\documentclass{article}

% Keep the package name that matches your local ICLR template files.
% If your template folder provides iclr2027_conference.sty, change 2026 to 2027.
\usepackage{iclr2026_conference,times}

\IfFileExists{math_commands.tex}{\input{math_commands.tex}}{}

\usepackage{booktabs}
\usepackage{multirow}
\usepackage{amsmath,amssymb}
\usepackage{graphicx}
\usepackage{hyperref}
\usepackage{tikz}
\usetikzlibrary{arrows.meta,calc,positioning,fit,backgrounds}
\usepackage{url}

\title{Are Neural Models Always Necessary For \\ Irregular Time Series Forecasting?}

% Anonymous submission. Replace after acceptance or when preparing a non-anonymous draft.
\author{Anonymous Authors}

\newcommand{\method}{AutoAnchor}
\newcommand{\family}{Anchor family}
\newcommand{\best}[1]{\textbf{#1}}
\newcommand{\second}[1]{\underline{#1}}

%\iclrfinalcopy

\begin{document}

\maketitle


\begin{abstract}
Irregular multivariate time-series (IMTS) forecasting is commonly assumed to require specialized neural architectures for their asynchronous nature. We audit this assumption by asking how much standard benchmark performance is explained by simple structure already present in the observed history. We introduce the Anchor family, a set of transparent, gradient-free diagnostic controls based on persistence, smoothing, sparse repetition, recurrence, and non-test empirical rule selection. Across global IMTS benchmarks and novel neural baselines under matched protocols, our Anchor variants achieve an optimal standard-window error in every dataset--metric setting, while AutoAnchor remains competitive across lookback--horizon sweeps and paired entity-level tests. Perturbation, thinning, bootstrap, and ablation analyses show that different datasets reward distinctive low-complexity historical structures, and that neural models regain advantages when history is insufficient or temporal configuration becomes demanding. These results suggest that current IMTS benchmarks can overstate the necessity of complex neural representations, thereby motivating stronger non-neural diagnostic controls first.
\end{abstract}


\section{Introduction}

Irregular multivariate time series (IMTS) forecasting is commonly treated as a
problem whose defining difficulties require specialized neural architectures.
Unlike regularly sampled sequences, IMTS data contain asynchronous observations,
variable-specific missingness, uneven time gaps, and partially observed future
targets. These properties have motivated architectures based on missingness-aware
recurrence~\citep{grud}, interpolation~\citep{shukla2019interp}, attention
over continuous time~\citep{mtan,primenet}, latent continuous
dynamics~\citep{neuralode2,latentODE,neuralCDE,neuralflows2,cru}, graph message
passing~\citep{raindrop2,gneuralflow,grafiti}, and adaptive
patching~\citep{tPatchGNN,apn}. The usual interpretation of leaderboard
improvements is that these models have learned useful irregular-time
representations.

This paper examines a narrower but important question: \emph{do today's standard
IMTS forecasting benchmarks require such representations to achieve strong
aggregate performance?} The distinction matters. A model can perform well on an
irregular benchmark because it reasons over asynchronous temporal structure, but
it can also perform well because the next observed value is often close to a
recent value, a smoothed level, a repeated sparse baseline, or a recurring local
pattern. In the latter case, a benchmark may reward anchoring to observed
history strongly enough that complex neural models appear necessary even when a
transparent deterministic rule explains much of the score. This study therefore frames the problem as a benchmark audit rather than as a new
state-of-the-art forecasting method. We introduce the \emph{Anchor family}, a
small ladder of deterministic diagnostic controls. NaiveAnchor repeats the most
recent value, ExpoAnchor applies horizon-aware exponential smoothing,
SparseAnchor probes repeated-value histories, ERMAnchor selects a candidate rule
by non-test empirical risk, and AutoAnchor augments this selection with generic
history-structure priors. Each rule is selected per variable using only
non-test data and is frozen before test evaluation. The family has no neural
parameters, no backpropagation, and no dataset- or variable-name lookup table;
its purpose is to estimate how much predictive signal is already present in the
observed history.

Across irregular time-series datasets, this diagnostic
control is surprisingly competitive with strong neural baselines under matched
benchmark protocols. The result should not be read as a claim that neural models
are generally unnecessary. Instead, it extends recent lessons from regular
forecasting benchmarks, where simple or lightweight models challenged the
necessity of increasingly complex Transformer variants~\citep{Zeng2023LTSFLinear,TFB2024,SimpleTM2025}.
Several established IMTS benchmarks contain enough persistence, smoothing,
sparse repetition, and recurrence for simple history-derived rules to match or
exceed learned models on important aggregate metrics. The scientific question is
therefore not only which architecture ranks first, but what predictive
information remains after these low-complexity structures have been controlled
for.

Our contributions are:
\begin{itemize}
    \item We formulate the Anchor family as transparent, gradient-free
    diagnostic controls for IMTS forecasting, designed to measure signal
    available without intricate learned representations.
    \item We conduct a matched benchmark audit against representative neural
    IMTS models using the same splits, preprocessing, masks, scaling, and
    evaluation metrics for fair comparison.
    \item We identify both regimes that reward anchoring and regimes where
    richer cross-variable temporal modeling remains vital, clarifying future IMTS benchmarking directions.
\end{itemize}


\section{Background and Motivation}

\begin{figure*}[t]
  \centering
  \includegraphics[width=\textwidth]{figures/general.jpg}
  \caption{General modeling pipeline for irregular multivariate time series by neural models \citep{cru, ContiFormer2023, grud, grafiti, mtan,seft, gneuralflow,tPatchGNN,apn, warpformer}. Overall, raw asynchronous observations are normalized and enriched with timing and missingness information, then embedded into a latent representation that captures temporal dynamics and cross-variable dependencies. This representation is finally decoded for downstream tasks.}
  \label{fig:imts-general-pipeline}
\end{figure*}

\paragraph{Learned representations for irregular time.}
Early IMTS methods made missingness and elapsed time explicit parts of the
model. GRU-D~\citep{grud} decays inputs and hidden states according to time
since observation, interpolation-prediction networks~\citep{shukla2019interp}
regularize sparse measurements before prediction, mTAN~\citep{mtan} uses
attention over continuous-time embeddings, and SeFT~\citep{seft} treats
irregular observations as a set rather than forcing them onto a regular grid.
Transformer-based approaches further model irregular temporal and cross-variable
interactions directly, including CoFormer~\citep{CoFormer2023}, while Rough
Transformers introduce signature-based continuous-time representations
~\citep{RoughTransformer2024}. A second line of work models hidden trajectories
in continuous time, beginning with Neural ODEs models~\citep{neuralode2, latentODE}, followed by neural controlled
differential equations~\citep{neuralCDE, neuralflows2,cru}. Graph- and patch-based methods  \citet{raindrop2,gneuralflow,grafiti,tPatchGNN,apn} add structure over
variables, observations, and temporal scales.
Recent work further expands this design space through channel-harmony
modeling~\citep{TimeCHEAT2025}, hierarchical patch graphs~\citep{HiPatch2025},
hypergraph dependency modeling~\citep{HyperIMTS2025}, and canonical
pre-alignment~\citep{KAFNet2026}. Figure~\ref{fig:imts-general-pipeline} summarizes the common neural pipeline.


\paragraph{Why stronger diagnostic controls are needed.}
The aforementioned IMTS architectures address real challenges in irregular forecasting, but benchmark performance alone does not reveal which modeling capabilities are actually necessary. Many datasets contain target values that are locally persistent, slowly varying, frequently repeated, or approximately recurrent, allowing expressive neural models to achieve strong aggregate scores partly by exploiting simple history-derived structure. This concern is not unique to IMTS. In long-sequence forecasting, increasingly specialized Transformer architectures were motivated by the need to capture complex temporal dependencies~\citep{informer}, yet later studies showed that direct multi-step linear models and other lightweight alternatives can remain surprisingly competitive under comparable benchmark protocols~\citep{Zeng2023LTSFLinear,SimpleTM2025}. Large-scale benchmarking further demonstrates that conclusions can change substantially when datasets, preprocessing, horizons, and evaluation procedures are standardized across methods~\citep{TFB2024}. We therefore frame IMTS benchmarking as a scientific control problem: before attributing performance gains to learned irregular-time representations, one should first establish how much of the benchmark signal can already be explained by transparent history-derived structure under the same splits, scaling, masking, horizons, and evaluation metrics.


\paragraph{Position of this work.}
We do not argue that recurrence, attention, latent dynamics, or graph learning are unnecessary for IMTS. Rather, their necessity should be demonstrated against controls that capture persistence, smoothing, sparse repetition, and recurrence under the same evaluation protocol. AutoAnchor serves as such a diagnostic: matching learned models indicates that benchmark scores alone do not establish the need for complex irregular-temporal or cross-variable representations, while its failures reveal regimes where neural modeling may provide genuine value. In addition, since recent methods such as APN~\citep{apn} already reduce computational cost, the central question is not efficiency, but whether increasing architectural complexity is empirically justified beyond strong non-neural controls, sparking pathways for customizing statistical models.




\section{The Anchor Family}

\begin{figure*}[t]
\centering
\small
\setlength{\tabcolsep}{5pt}
\renewcommand{\arraystretch}{1.35}

% ===== Outer border =====
\setlength{\fboxsep}{8pt}
\fbox{%
\begin{minipage}{0.96\textwidth}
\centering

% ===== Shared anchor class =====
\setlength{\fboxsep}{3pt}
\fbox{%
\begin{minipage}{0.961\textwidth}
\centering
\vspace{0.25em}
\textbf{Shared anchor class}
\[
\mathcal A=\left\{a:
 f_a(\mathcal H_v,t^*)=\sum_{\ell=1}^{J_a}w_{a,\ell}
 g_{a,\ell}(\mathcal H_v,t^*),\quad
 w_{a,\ell}\ge 0,\ \sum_{\ell}w_{a,\ell}=1
\right\}.
\]
\vspace{-0.8em}
% {\footnotesize Primitive rules $g$ include last value, local level summaries,
% mode, EMA, bounded trend, and phase-nearest lookup.}
\vspace{0.35em}
\end{minipage}%
}

\vspace{0.55em}

% ===== Anchor variants =====
\begin{tabular}{p{0.15\textwidth}|p{0.47\textwidth}|p{0.28\textwidth}}
\toprule
\textbf{Variant} & \textbf{Calibration rule for variable $v$} & \textbf{Frozen specification} \\
\midrule
\textbf{NaiveAnchor}
& Prescribe the persistence candidate.
& $\theta_v=(\mathrm{Last},1)$ \\
\cmidrule(lr){1-3}

\textbf{ExpoAnchor}
& Prescribe an EMA candidate with horizon-aware smoothing,
$\alpha=\mathrm{clip}\{\mathrm{round}(10\sqrt{K/L})/10,0.1,0.9\}$.
& $\theta_v=(\mathrm{EMA}_{\alpha},1)$ \\
\cmidrule(lr){1-3}

\textbf{SparseAnchor}
& Apply a sparse-history diagnostic $s_v=(p_v,u_v)$, where $p_v$ is the
dominant-value frequency and $u_v$ is the rounded-value cardinality; choose a
mode--phase candidate when the diagnostic activates.
& $\theta_v=(a_v^{\mathrm{sp}},1)$ \\
\cmidrule(lr){1-3}

\textbf{ERMAnchor}
& Select by non-test empirical risk:
$\displaystyle
(a_v,\beta_v)=
\arg\min_{a\in\mathcal A,\,\beta\in B_a}
\widehat R_v(a,\beta)$.
& $\theta_v=(a_v^{\mathrm{ERM}},\beta_v)$ \\
\cmidrule(lr){1-3}

\textbf{AutoAnchor}
& Use an ordered structural prior $\pi_v(\mathcal H)$ when it fires; otherwise
retain the ERM solution:
$\displaystyle
\theta_v=
\begin{cases}
(\pi_v(\mathcal H),1),&\pi_v\neq\varnothing,\\
\theta_v^{\mathrm{ERM}},&\text{otherwise.}
\end{cases}$
& $\theta_v=(a_v^{\mathrm{prior}},1)$ or
$\theta_v^{\mathrm{ERM}}$ \\
\bottomrule
\end{tabular}

\vspace{0.55em}

% ===== Shared test-time operator =====
\fbox{%
\begin{minipage}{0.961\textwidth}
\centering
\textbf{Shared test-time operator}\qquad
$\displaystyle
\hat y_{b,k,v}
=
P_{b,v}\!\left(
\beta_v f_{a_v}(\mathcal H_{b,v},t^*_{b,k,v})
\right),
$
\quad where $P_{b,v}$ denotes the Anchor output projection after inverse scaling.
\end{minipage}%
}

\end{minipage}%
}
% ===== End outer border =====

\caption{\textbf{Anchor Family architecture.}
All variants share the same deterministic candidate class $\mathcal A$ and
differ only in how the per-variable specification
$\theta_v=(a_v,\beta_v)$ is calibrated. The frozen specification is then
applied through a common test-time forecast operator.}
\label{fig:anchor-family-architecture}
\end{figure*}

The Anchor family is a diagnostic ladder of deterministic forecasting controls.
All variants operate on the same candidate class $\mathcal A$ in
Figure~\ref{fig:anchor-family-architecture}. Each candidate combines a small
number of primitive history rules, such as last-value persistence, recent-level
summaries, trimmed means, mode, exponential smoothing, bounded local trend, and
phase-nearest lookup. A candidate forecast shows:
\begin{equation}
 f_a(\mathcal H,t^*)=\sum_{\ell=1}^{J_a}w_{a,\ell}
 g_{a,\ell}(\mathcal H,t^*),\qquad
 w_{a,\ell}\geq0,\quad \sum_{\ell}w_{a,\ell}=1.
 \label{eq:anchor-candidate}
\end{equation}
The library is fixed before evaluation and shared across datasets and
variables. It is not intended to be an optimal hand-engineered forecaster; it is
intended to span common non-neural explanations for benchmark success. For each variable $v$, the adaptive variants select a candidate $a$ and an
optional scalar coefficient $\beta$ using only train and validation targets:
\begin{equation}
 \widehat R_v(a,\beta)=
 \frac{\sum_{b,k}m_{b,k,v}
 \left(\widetilde y_{b,k,v}-\widetilde f_{a,\beta}
 (\mathcal aH_{b,v},t^*_{b,k,v})\right)^2}
 {\sum_{b,k}m_{b,k,v}},\qquad b\in\mathcal D_{\mathrm{cal}}.
 \label{eq:anchor-risk}
\end{equation}
Accordingly, $m$ is the target mask and tildes denote benchmark-scaled quantities after
the same deterministic Anchor output operator used at test time. The resulting specification
$\theta_v=(a_v,\beta_v)$ is frozen before test evaluation. There is no
sample-wise test-time search over the library. The five variants differ only in how $\theta_v$ is chosen. NaiveAnchor fixes the
last-value rule, measuring pure persistence. ExpoAnchor fixes an EMA rule whose
smoothing depends on the ratio between forecast horizon and lookback length.
SparseAnchor uses a simple repeated-value diagnostic to choose between mode,
mode--phase, and last-value rules. ERMAnchor chooses the candidate and scale
coefficient with the lowest non-test risk in Equation~\ref{eq:anchor-risk}.
AutoAnchor starts from ERMAnchor, then applies an ordered structural prior when
calibration histories show strong repetition, long-window short-horizon
smoothing, or phase-like recurrence:
\begin{equation}
 \theta_v^{\mathrm{auto}}=
 \begin{cases}
 (a_v^{\mathrm P},1),&\text{if a structural diagnostic activates},\\
 \theta_v^{\mathrm{ERM}},&\text{otherwise}.
 \end{cases}
 \label{eq:auto-summary}
\end{equation}
The prior is dataset-agnostic and does not read variable names. Its thresholds
are disclosed exploratory choices rather than universal constants, so we
interpret AutoAnchor together with the full anchor ladder and ablations.
Candidate rules are fixed in the released
implementation.


\section{Benchmark Audit Protocol}

\paragraph{Task and notation.}
For each sample $b$, the observed history is an irregular set of timestamped
measurements. For a variable $v$, we write
\begin{equation}
    \mathcal H_{b,v}=\{(t_{b,v,j},x_{b,v,j})\}_{j=1}^{n_{b,v}},
    \qquad t_{b,v,1}<\cdots<t_{b,v,n_{b,v}}.
\end{equation}
The goal is to predict observed future targets $y_{b,k,v}$ at query times
$t^*_{b,k,v}$. Since future targets are partially observed, all
metrics use the target mask $m_{b,k,v}\in\{0,1\}$ supplied by the benchmark.

\paragraph{Datasets.}
We audit four widely used IMTS forecasting benchmarks. PhysioNet Challenge 2012
(P12) contains ICU time-varying variables~\citep{Silva2012PhysioNet2}; MIMIC
contains de-identified critical-care trajectories~\citep{Johnson2016MIMIC};
HumanActivity contains wearable-sensor trajectories for human motion
classification/forecasting settings~\citep{apn}; and USHCN
contains irregular climate observations from United States weather
stations~\citep{Menne2016USHCN}. These datasets differ substantially in
sampling density, target availability, variable volatility, and temporal scale,
which makes them useful for diagnosing when anchoring is sufficient and when it
is not.

\paragraph{Baseline Models.}
We compare against representative IMTS models spanning missingness-aware
recurrence, set encoders, attention, latent dynamics, graph modeling, and
adaptive patching: PrimeNet~\citep{primenet}, NeuralFlows~\citep{neuralflows2},
CRU~\citep{cru}, mTAN~\citep{mtan}, SeFT~\citep{seft}, GNeuralFlow~\citep{gneuralflow},
GRU-D~\citep{grud}, Raindrop~\citep{raindrop2}, Warpformer~\citep{warpformer},
tPatchGNN~\citep{tPatchGNN}, GraFITi~\citep{grafiti}, and APN~\citep{apn}. We
follow \citet{apn}'s benchmark protocol as standards for splits, sequence lengths,
prediction lengths, masks, scaling, and global metric definitions. To ensure a fair comparison, all neural
baselines use the released hyperparameter
settings in original papers while
holding the dataset splits, preprocessing, prediction horizons, and evaluation
protocol fixed across methods; no test-label tuning or sample-wise test routing is used.
Results are reported over multiple seeds (2024-2026). We also avoid evaluation shortcuts such as the ``Drop Last'' trick
discussed in recent forecasting benchmark audits~\citep{TFB2024,TimeRecipe2026}.
Anchor methods are deterministic diagnostic controls and therefore do not require
stochastic training, optimizers, epochs, or learned neural parameters.

\paragraph{Evaluation Metrics.}
Let $s_v(\cdot)$ be the training-split scaling transform for variable $v$. We
evaluate predictions in scaled space:
\begin{equation}
    \widetilde y_{b,k,v}=s_v(y_{b,k,v}),\qquad
    \widetilde{\hat y}_{b,k,v}=s_v(\hat y_{b,k,v}).
\end{equation}
The global masked metrics are
\begin{equation}
    \mathrm{MAE}=\frac{\sum_{b,k,v}m_{b,k,v}
    |\widetilde{\hat y}_{b,k,v}-\widetilde y_{b,k,v}|}
    {\sum_{b,k,v}m_{b,k,v}},
    \qquad
    \mathrm{MSE}=\frac{\sum_{b,k,v}m_{b,k,v}
    (\widetilde{\hat y}_{b,k,v}-\widetilde y_{b,k,v})^2}
    {\sum_{b,k,v}m_{b,k,v}}.
\end{equation}
We use global masked metrics because they are the benchmark convention, but we
also report paired entity-level tests and bootstrap intervals to avoid relying
only on aggregate leaderboard differences.

\paragraph{Audit design.}
The audit asks whether neural gains survive strong non-neural controls under the
same protocol. Table~\ref{tab:main-results} establishes the headline
benchmark comparison at each dataset's standard window, asking whether learned
architectures retain an aggregate advantage once strong deterministic controls
are included. Figure~\ref{fig:cross-baseline-lookback} then varies lookback and
prediction horizon to test whether the conclusion is tied to a single favorable
window. Table~\ref{tab:paired-neural-tests} replaces leaderboard averages with
paired entity-level sign-flip tests, measuring whether observed differences are
consistent across patients, stations, or activity windows rather than driven by
a few samples. Table~\ref{tab:cross-model-temporal}
and Figure~\ref{fig:cross-model-sparsity} intervene on timestamp order and available
history to identify whether performance depends on irregular temporal geometry
or mostly on recent observed values. Finally, Table~\ref{tab:bootstrap-ci}
quantifies uncertainty for Anchor-family metrics, while
Table~\ref{tab:candidate-removal} attributes performance to specific candidate
classes.

% \newpage
\section{Main Results}

% The experiments answer three questions. First, can simple history-derived
% controls match neural IMTS models at the standard benchmark settings? Second,
% does that conclusion survive changes in lookback/horizon and paired entity-level
% significance tests? Third, which observed-history properties does performance depend on, and
% where do Anchor controls stop being sufficient?

\begin{table*}[t]
\caption{Comparison at the standard benchmark window of each dataset. All
methods use a common data and evaluation protocol. Neural methods report mean
$\pm$ standard deviation over three runs (seed 2024-2026); Anchor methods are deterministic.
Lower is better. \textbf{Bold} marks the best value in each metric column and
\underline{underlining} marks the second best.}
\label{tab:main-results}
\centering
\scriptsize
\resizebox{\textwidth}{!}{
\begin{tabular}{l|c|c|c|c|c|c|c|c}
\toprule
Method
& \multicolumn{2}{|c|}{HumanActivity}
& \multicolumn{2}{c|}{USHCN}
& \multicolumn{2}{c|}{P12}
& \multicolumn{2}{c}{MIMIC} \\
\cmidrule(lr){2-3}
\cmidrule(lr){4-5}
\cmidrule(lr){6-7}
\cmidrule(lr){8-9}
& MSE & MAE & MSE & MAE & MSE & MAE & MSE & MAE \\
\midrule
PrimeNet & 4.2507$\pm$0.0041 & 1.7018$\pm$0.0011 & 0.4930$\pm$0.0015 & 0.4954$\pm$0.0018 & 0.7953$\pm$0.0000 & 0.6859$\pm$0.0001 & 0.9073$\pm$0.0001 & 0.6614$\pm$0.0001 \\
NeuralFlows & 0.1722$\pm$0.0090 & 0.3150$\pm$0.0094 & 0.2087$\pm$0.0258 & 0.3157$\pm$0.0187 & 0.4056$\pm$0.0033 & 0.4466$\pm$0.0027 & 0.6085$\pm$0.0101 & 0.5306$\pm$0.0066 \\
CRU & 0.1387$\pm$0.0073 & 0.2607$\pm$0.0092 & 0.2168$\pm$0.0162 & 0.3180$\pm$0.0248 & 0.6179$\pm$0.0045 & 0.5778$\pm$0.0031 & 0.5895$\pm$0.0092 & 0.5151$\pm$0.0048 \\
mTAN & 0.0993$\pm$0.0026 & 0.2219$\pm$0.0047 & 0.5561$\pm$0.2020 & 0.5015$\pm$0.0968 & 0.3809$\pm$0.0043 & 0.4291$\pm$0.0035 & 0.9408$\pm$0.1126 & 0.6755$\pm$0.0459 \\
SeFT & 1.3786$\pm$0.0024 & 0.9762$\pm$0.0007 & 0.3345$\pm$0.0022 & 0.4083$\pm$0.0084 & 0.7721$\pm$0.0021 & 0.6760$\pm$0.0029 & 0.9230$\pm$0.0015 & 0.6628$\pm$0.0008 \\
GNeuralFlow & 0.3936$\pm$0.1585 & 0.4541$\pm$0.0841 & 0.2205$\pm$0.0421 & 0.3286$\pm$0.0412 & 0.8207$\pm$0.0310 & 0.6759$\pm$0.0100 & 0.8957$\pm$0.0209 & 0.6450$\pm$0.0072 \\
GRU-D & 0.1893$\pm$0.0627 & 0.3253$\pm$0.0485 & 0.2097$\pm$0.0493 & 0.3045$\pm$0.0305 & 0.3419$\pm$0.0029 & 0.3992$\pm$0.0011 & 0.4759$\pm$0.0100 & 0.4526$\pm$0.0055 \\
Raindrop & 0.0916$\pm$0.0072 & 0.2114$\pm$0.0072 & 0.2035$\pm$0.0336 & 0.3029$\pm$0.0264 & 0.3478$\pm$0.0019 & 0.4044$\pm$0.0020 & 0.6754$\pm$0.1829 & 0.5444$\pm$0.0868 \\
Warpformer & 0.0449$\pm$0.0010 & 0.1228$\pm$0.0018 & 0.1888$\pm$0.0598 & 0.2939$\pm$0.0591 & \second{0.3056$\pm$0.0011} & 0.3661$\pm$0.0016 & 0.4302$\pm$0.0035 & 0.4025$\pm$0.0014 \\
tPatchGNN & 0.0443$\pm$0.0009 & 0.1247$\pm$0.0031 & 0.1885$\pm$0.0403 & 0.3084$\pm$0.0479 & 0.3133$\pm$0.0053 & 0.3697$\pm$0.0049 & 0.4431$\pm$0.0115 & 0.4077$\pm$0.0088 \\
GraFITi & 0.0437$\pm$0.0005 & 0.1221$\pm$0.0017 & 0.1691$\pm$0.0093 & 0.2777$\pm$0.0248 & 0.3075$\pm$0.0015 & \second{0.3637$\pm$0.0036} & 0.4359$\pm$0.0455 & 0.4142$\pm$0.0297 \\
APN & \second{0.0421$\pm$0.0001} & 0.1159$\pm$0.0006 & \second{0.1590$\pm$0.0137} & \second{0.2611$\pm$0.0167} & 0.3093$\pm$0.0011 & 0.3650$\pm$0.0026 & 0.4292$\pm$0.0027 & 0.4016$\pm$0.0016 \\
\midrule
NaiveAnchor & 0.0605 & 0.1225 & 0.3640 & 0.3291 & 0.3992 & 0.3951 & 0.4561 & 0.3708 \\
ExpoAnchor & \best{0.0420} & \second{0.1148} & 0.5923 & 0.5117 & 0.3286 & 0.3753 & 0.6431 & 0.4974 \\
SparseAnchor & 0.0605 & 0.1225 & 0.2661 & 0.2812 & 0.3992 & 0.3951 & 0.4560 & 0.3709 \\
ERMAnchor & 0.0428 & \best{0.1131} & 0.1662 & 0.2635 & \best{0.3020} & \best{0.3585} & \best{0.3782} & \second{0.3658} \\
\method{} & \best{0.0420} & \second{0.1148} & \best{0.1565} & \best{0.2187} & \best{0.3020} & \best{0.3585} & \second{0.3785} & \best{0.3657} \\
\bottomrule
\end{tabular}
}
\end{table*}


\begin{figure*}[t]
\centering
\begin{minipage}[t]{0.49\textwidth}
  \centering
  \includegraphics[width=\linewidth]{figures/P12_lookback_horizon.png}
  \centerline{\footnotesize (a) P12}
\end{minipage}\hfill
\begin{minipage}[t]{0.49\textwidth}
  \centering
  \includegraphics[width=\linewidth]{figures/MIMIC_lookback_horizon.png}
  \centerline{\footnotesize (b) MIMIC}
\end{minipage}
\vspace{0.6em}

\begin{minipage}[t]{0.49\textwidth}
  \centering
  \includegraphics[width=\linewidth]{figures/HumanActivity_lookback_horizon.png}
  \centerline{\footnotesize (c) HumanActivity}
\end{minipage}\hfill
\begin{minipage}[t]{0.49\textwidth}
  \centering
  \includegraphics[width=\linewidth]{figures/USHCN_lookback_horizon.png}
  \centerline{\footnotesize (d) USHCN}
\end{minipage}
\caption{Lookback--horizon sweep for \method{} and the Top 3 neural models.  The plots show whether the Anchor conclusion persists across configurations and where others regain an advantage.}
\label{fig:cross-baseline-lookback}
\end{figure*}

\begin{table*}[t]
\caption{Paired entity-level tests comparing \method{} with neural baselines at the standard benchmark windows. Entries are paired entity-level averages used for the test, rather than the global aggregates in Table~\ref{tab:main-results}. All paired comparisons use the common valid target-coordinate set after aligning the rerun exports. Deltas are \method{} minus neural error, so negative values favor \method{}. Reported $p$-values are unadjusted paired entity-level sign-flip tests; the text summarizes Holm-corrected counts within each metric family.}
\label{tab:paired-neural-tests}
\centering
\scriptsize
\resizebox{\textwidth}{!}{
\begin{tabular}{l|l|c|c|c|c|c|c|c|c}
\toprule
Dataset & Model & Auto MAE & Neural MAE & $\Delta$ MAE & $p$ & Auto MSE & Neural MSE & $\Delta$ MSE & $p$ \\
\midrule
\multirow{3}{*}{P12} & APN & 0.3585 & 0.3649 &\best{ -0.0063 }& $<10^{-4}$ & 0.3020 & 0.3032 &\best{ -0.0012 }& 0.4953 \\
& GraFITi & 0.3585 &\best{ 0.3581 }& +0.0004 & 0.6687 & 0.3020 &\best{ 0.2953 }& +0.0067 & 0.0002 \\
& tPatchGNN & 0.3585 &\second{ 0.3628 }&\second{ -0.0043 }& $<10^{-4}$ & 0.3020 &\second{ 0.3006 }& +0.0014 & 0.9509 \\
\midrule
\multirow{3}{*}{MIMIC} & APN & 0.3657 &\second{ 0.3737 }&\second{ -0.0080 }& $<10^{-4}$ & 0.3785 &\best{ 0.3931 }&\second{ -0.0147 }& 0.2684 \\
& GraFITi & 0.3657 & 0.4150 &\best{ -0.0493 }& $<10^{-4}$ & 0.3785 & 0.4365 &\best{ -0.0580 }& $<10^{-4}$ \\
& tPatchGNN & 0.3657 &\best{ 0.3676 }& -0.0019 & 0.0020 & 0.3785 &\second{ 0.3948 }&\second{ -0.0163 }& 0.1895 \\
\midrule
\multirow{3}{*}{USHCN} & APN & 0.2187 & 0.2844 &\best{ -0.0657 }& $<10^{-4}$ & 0.1565 &\second{ 0.1611 }&\second{ -0.0046 }& 0.8642 \\
& GraFITi & 0.2187 &\second{ 0.2730 }&\second{ -0.0543 }& $<10^{-4}$ & 0.1565 & 0.1669 &\best{ -0.0104 }& 0.6827 \\
& tPatchGNN & 0.2187 &\best{ 0.2668 }& -0.0481 & 0.0011 & 0.1565 &\best{ 0.1596 }& -0.0031 & 0.9600 \\
\midrule
\multirow{3}{*}{HumanActivity} & APN & 0.1148 &\best{ 0.1156 }& -0.0009 & 0.1383 & 0.0420 &\best{ 0.0421 }& -0.0001 & 0.8420 \\
& GraFITi & 0.1148 &\second{ 0.1226 }&\second{ -0.0078 }& $<10^{-4}$ & 0.0420 &\second{ 0.0439 }&\second{ -0.0019 }& 0.0085 \\
& tPatchGNN & 0.1148 & 0.1323 &\best{ -0.0175 }& $<10^{-4}$ & 0.0420 & 0.0470 &\best{ -0.0051 }& $<10^{-4}$ \\
\bottomrule
\end{tabular}}
\end{table*}

\begin{table*}[t]
\caption{Cross-model temporal perturbations at the standard benchmark windows. Original MSE is the unperturbed error; each $\Delta$ is the mean MSE change under the perturbation. Positive values denote degradation and measure sensitivity, not ranking quality.}
\label{tab:cross-model-temporal}
\centering
\scriptsize
\resizebox{\textwidth}{!}{
\begin{tabular}{l|l|c|c|c|c|c}
\toprule
Dataset & Model & Orig. MSE & Shuffle $\Delta$ & Swap $\Delta$ & Last-only $\Delta$ & Gaps-only $\Delta$ \\
\midrule
\multirow{4}{*}{P12} & \method{} &\best{ 0.3020 }& +0.0007 & +0.2240 & +0.0575 & +0.4997 \\
& APN & 0.3093$\pm$0.0011 & +0.1467 & +0.2023 & +0.0558 & +0.4873 \\
& GraFITi &\second{ 0.3075$\pm$0.0015 }& +0.1471 & +0.1900 & +0.0464 & +0.4906 \\
& tPatchGNN & 0.3133$\pm$0.0053 & +0.4549 & +0.1826 & +0.0416 & +0.4827 \\
\midrule
\multirow{4}{*}{MIMIC} & \method{} &\best{ 0.3785 }& +0.0102 & +0.3756 & +0.0076 & +0.6575 \\
& APN &\second{ 0.4292$\pm$0.0027 }& +0.1843 & +0.2344 & +0.0454 & +0.5408 \\
& GraFITi & 0.4359$\pm$0.0455 & +0.1802 & +0.2359 & +0.0679 & +0.5348 \\
& tPatchGNN & 0.4431$\pm$0.0115 & +0.1809 & +0.2372 & +0.0444 & +0.5170 \\
\midrule
\multirow{4}{*}{USHCN} & \method{} &\best{ 0.1565 }& +0.3125 & +0.6718 & +0.2075 & +0.3362 \\
& APN &\second{ 0.1590$\pm$0.0137 }& +0.0023 & +0.0045 & +0.1401 & +0.0804 \\
& GraFITi & 0.1691$\pm$0.0093 & -0.0015 & -0.0019 & +0.0473 & +0.1010 \\
& tPatchGNN & 0.1885$\pm$0.0403 & -0.0195 & +0.3605 & +0.0595 & +0.1267 \\
\midrule
\multirow{4}{*}{HumanActivity} & \method{} &\best{ 0.0420 }& +0.0000 & +0.0611 & +0.0185 & +4.2071 \\
& APN &\second{ 0.0421$\pm$0.0001 }& +0.0349 & +0.0581 & +0.0166 & +4.0887 \\
& GraFITi & 0.0437$\pm$0.0005 & +0.0347 & +0.0488 & +0.1217 & +3.5537 \\
& tPatchGNN & 0.0443$\pm$0.0009 & +0.3731 & +0.0554 & +0.2989 & +3.7591 \\
\bottomrule
\end{tabular}}
\end{table*}


\begin{figure*}[t]
\centering

\makebox[\textwidth][c]{%
\begin{minipage}[t]{0.35\textwidth}
  \centering
  \includegraphics[width=\linewidth]{figures/HumanActivity_history_thinning.png}
  \centerline{\footnotesize (a) HumanActivity}
\end{minipage}
\hspace{0.04\textwidth}
\begin{minipage}[t]{0.35\textwidth}
  \centering
  \includegraphics[width=\linewidth]{figures/MIMIC_history_thinning.png}
  \centerline{\footnotesize (b) MIMIC}
\end{minipage}
}

\vspace{0.6em}

\makebox[\textwidth][c]{%
\begin{minipage}[t]{0.35\textwidth}
  \centering
  \includegraphics[width=\linewidth]{figures/P12_history_thinning.png}
  \centerline{\footnotesize (c) P12}
\end{minipage}
\hspace{0.04\textwidth}
\begin{minipage}[t]{0.35\textwidth}
  \centering
  \includegraphics[width=\linewidth]{figures/USHCN_history_thinning.png}
  \centerline{\footnotesize (d) USHCN}
\end{minipage}
}

\caption{Cross-model history-thinning tests at the standard benchmark windows. Each panel reports global scaled MSE after retaining a seeded random fraction of observed history entries at evaluation time while preserving remaining timestamps and values.}
\label{fig:cross-model-sparsity}
\end{figure*}


\begin{table*}[t]
\caption{Entity-bootstrap 95\% confidence intervals over patients, stations, or
activity windows for all Anchor-family global metrics. MSE red. is the relative
MSE reduction against NaiveAnchor within each dataset; negative values indicate
higher MSE than NaiveAnchor.}
\label{tab:bootstrap-ci}
\centering
\scriptsize
\resizebox{\textwidth}{!}{
\begin{tabular}{l|l|c|c|c}
\toprule
Dataset & Method & MAE (95\% CI) & MSE (95\% CI) & MSE red. (\%) \\
\midrule
\multirow{5}{*}{P12} & NaiveAnchor & 0.3951 [0.3834, 0.4076] & 0.3992 [0.3741, 0.4273] & 0.0\% \\
& ExpoAnchor & 0.3753 [0.3653, 0.3856] & 0.3286 [0.3078, 0.3521] & 17.7\% \\
& SparseAnchor & 0.3951 [0.3834, 0.4076] & 0.3992 [0.3741, 0.4273] & 0.0\% \\
& ERMAnchor & 0.3585 [0.3488, 0.3686] & 0.3020 [0.2820, 0.3245] & 24.3\% \\
& \method{} & 0.3585 [0.3488, 0.3686] & 0.3020 [0.2820, 0.3245] & 24.3\% \\
\midrule
\multirow{5}{*}{MIMIC} & NaiveAnchor & 0.3708 [0.3569, 0.3853] & 0.4561 [0.4136, 0.5034] & 0.0\% \\
& ExpoAnchor & 0.4974 [0.4809, 0.5139] & 0.6431 [0.5850, 0.7080] & -41.0\% \\
& SparseAnchor & 0.3709 [0.3571, 0.3855] & 0.4560 [0.4137, 0.5033] & 0.0\% \\
& ERMAnchor & 0.3658 [0.3537, 0.3785] & 0.3782 [0.3396, 0.4242] & 17.1\% \\
& \method{} & 0.3657 [0.3536, 0.3784] & 0.3785 [0.3397, 0.4242] & 17.0\% \\
\midrule
\multirow{5}{*}{USHCN} & NaiveAnchor & 0.3291 [0.2787, 0.3837] & 0.3640 [0.2640, 0.4791] & 0.0\% \\
& ExpoAnchor & 0.5117 [0.4519, 0.5751] & 0.5923 [0.4955, 0.6994] & -62.8\% \\
& SparseAnchor & 0.2812 [0.2382, 0.3266] & 0.2661 [0.2029, 0.3367] & 26.9\% \\
& ERMAnchor & 0.2634 [0.2315, 0.2963] & 0.1662 [0.1289, 0.2074] & 54.3\% \\
& \method{} & 0.2187 [0.1868, 0.2528] & 0.1565 [0.1209, 0.1958] & 57.0\% \\
\midrule
\multirow{5}{*}{HumanActivity} & NaiveAnchor & 0.1225 [0.1116, 0.1341] & 0.0605 [0.0480, 0.0746] & 0.0\% \\
& ExpoAnchor & 0.1148 [0.1054, 0.1248] & 0.0420 [0.0342, 0.0504] & 30.6\% \\
& SparseAnchor & 0.1225 [0.1116, 0.1341] & 0.0605 [0.0480, 0.0746] & 0.0\% \\
& ERMAnchor & 0.1131 [0.1039, 0.1230] & 0.0428 [0.0346, 0.0519] & 29.2\% \\
& \method{} & 0.1148 [0.1054, 0.1248] & 0.0420 [0.0342, 0.0504] & 30.6\% \\
\bottomrule
\end{tabular}}
\end{table*}


\paragraph{Standard-window benchmark audit.}
Table~\ref{tab:main-results} shows that the Anchor family is more than a sanity-check baseline at standard benchmark windows. An Anchor variant achieves the best or tied-best result in all eight dataset--metric cells, with \method{} doing so in six. The effect sizes, however, vary substantially. Relative to the strongest neural baseline, gains are modest on HumanActivity MSE (0.24\%) and P12 MSE/MAE (1.2--1.4\%), but larger on USHCN MAE (16.2\%) and MIMIC MSE/MAE (11.9\% and 8.9\%). The key result is therefore not that every improvement is equally large, but that competitive benchmark error can often be matched without learned representations. The Anchor ladder further shows that no single heuristic explains all datasets: HumanActivity is largely captured by horizon-aware smoothing, P12 and MIMIC by empirical selection over history-derived rules, and USHCN by sparse or phase structure. Different benchmarks thus reward different low-complexity structures already present in the observed history.


\paragraph{Lookback and horizon robustness.}
Figure~\ref{fig:cross-baseline-lookback} tests whether the standard-window
finding is tied to a favorable configuration. Across the 19 tested windows,
\method{} obtains the lowest MSE in 13 and the lowest MAE in 16, including all
six MIMIC configurations and four of five P12 configurations for both metrics.
The controlled window changes are also informative. On USHCN, moving from
$(L,H)=(100,3)$ to $(150,3)$ reduces \method{} MSE from 1.1591 to 0.1565,
consistent with longer historical context being important for the
phase-sensitive behavior identified by the candidate analysis.
At fixed $L=150$, increasing the horizon from $H=3$ to $H=7$ lets neural models
regain the MSE advantage, while \method{} retains the MAE advantage. These
exceptions, together with the small APN MSE advantages on some HumanActivity
windows, identify regimes where learned temporal aggregation can yield measurable
gains beyond frozen controls.

\paragraph{Paired comparison with neural models.}
Table~\ref{tab:paired-neural-tests} replaces aggregate leaderboard comparisons
with paired tests on matched entities. After Holm correction within each metric
family, \method{} retains significant MAE improvements in 10 of the 12 neural
comparisons. MSE remains more mixed: three of 12 differences remain significant,
two favoring \method{} and one favoring GraFITi on P12. This distinction is
central to the benchmark-audit framing. Paired analysis shows that several small
differences between neural models and \method{} are not statistically
distinguishable at the entity level, while a smaller subset exhibits clear
advantages in either direction.

\paragraph{Temporal perturbations.}
Table~\ref{tab:cross-model-temporal} should be read as a sensitivity analysis,
not as a leaderboard: a smaller $\Delta$ means lower sensitivity to the
perturbed information, not necessarily a better model. Small negative deltas can
arise from finite-sample variation or perturbation-induced changes in model
behavior and should not be interpreted as evidence that removing temporal
information is beneficial. On P12, MIMIC, and HumanActivity, timestamp
shuffling changes \method{} MSE by only about 0.001,
0.010, and 0.000, respectively, even though several neural models deteriorate
substantially under the same intervention. This supports a benchmark-level
conclusion: exact timestamp geometry may be used by neural models, but it is not
necessary for competitive aggregate performance on these three settings. USHCN
is the counterexample. For \method{}, timestamp shuffling increases MSE by
0.313 and swapping history halves increases it by 0.672, showing that
\method{}'s strong USHCN performance depends substantially on temporal
arrangement.

\paragraph{History thinning.}
Figure~\ref{fig:cross-model-sparsity} reports a seeded random thinning study that retains a fixed fraction of
observed history entries while preserving the timestamps and values that remain.
For this study, thinning is applied only at evaluation time: neural checkpoints
are trained on original histories, and Anchor specifications are calibrated on
original histories before the thinned test histories are applied.
The dataset-level patterns differ sharply. On P12, \method{} remains best at
10\% history, but only narrowly ahead of GraFITi (0.5224 versus 0.5279). On
MIMIC, the advantage is much clearer at the same retention level (0.7253 versus
0.8375 for the best neural comparator). USHCN reverses under aggressive
thinning, with neural models substantially ahead at 10\% history, and
HumanActivity shifts from an AutoAnchor advantage at 10\% to small APN
advantages at intermediate fractions. These reversals are scientifically
informative: they identify settings where neural models exhibit greater
test-time robustness than frozen history-based controls under severe history
loss. Additionally, the relative MSE reductions, ranging from 17\% to 57\%, over NaiveAnchor show that the Anchor family provides substantial gains beyond pure persistence, indicating that its performance comes from structured history-based selection rather than last-value copying alone.

\paragraph{Uncertainty and candidate attribution.}
Table~\ref{tab:bootstrap-ci} reports entity-bootstrap intervals for the full
Anchor ladder. The aggregate estimates are stable under resampling: \method{}
and NaiveAnchor have clearly separated MSE intervals on P12 and USHCN, while
HumanActivity shows a large MSE reduction with slight interval overlap. The full
ladder is also informative because some fixed rules are neutral or harmful, such
as ExpoAnchor on MIMIC and USHCN, whereas ERMAnchor and \method{} recover strong
reductions over persistence. These intervals summarize sampling uncertainty
rather than formal paired significance. Table~\ref{tab:candidate-removal} then
attributes performance to candidate classes. P12 and MIMIC are largely matched
by ERM selection from the shared library; USHCN benefits strongly from
structural sparse/phase components; and HumanActivity is mostly matched by the
horizon-aware EMA rule. The key benchmark-audit conclusion is that different
datasets reward different low-complexity structures, yet all four standard
benchmark settings are substantially matched by transparent history-derived
controls.

\paragraph{Failure modes.}
The real-data reversals already define the boundary of the claim. Neural models
regain advantages on USHCN when the lookback is too short or the horizon is
extended, and on HumanActivity under several window and thinning regimes. These
failures are expected for a variable-wise Anchor family: it does not explicitly
learn latent state, cross-variable dependencies, or anticipatory regime
dynamics. The paper therefore does not argue that neural IMTS modeling is
unnecessary; it argues that its advantages
should be demonstrated on settings where this additional representational
capacity is required and measurable.

% \paragraph{Synthesis.}
% Together, the experiments show a coherent pattern. Standard IMTS benchmark
% performance can often be matched by low-complexity history controls; this
% finding persists across many windows and paired entity-level comparisons; and
% perturbations reveal that different datasets reward different forms of
% persistence, smoothing, sparsity, and recurrence. At the same time, the same
% audit exposes where anchors fail and neural models regain an advantage. Thus,
% the benchmark question is not simply whether a neural architecture achieves
% lower error, but whether its gain remains after controlling for low-complexity
% historical structure.

\section{Discussion}

\paragraph{What the audit establishes.}
Our results show that strong IMTS benchmark performance does not necessarily imply that complex irregular-time representations are required. Across the standard settings, transparent history-derived controls match or exceed specialized neural models, and paired tests show that many remaining differences are small or statistically indistinguishable. The key implication is therefore diagnostic: benchmark accuracy should not be treated as evidence of representational necessity unless gains remain after controlling for low-complexity historical structure.

\paragraph{Why the result is not a trivial-baseline artifact.}
No single heuristic explains all four datasets. HumanActivity is largely captured by horizon-aware smoothing, P12 and MIMIC by empirical selection over simple history-based rules, and USHCN by longer-context sparse/phase structure. This heterogeneity strengthens the audit: different benchmarks reward different forms of low-complexity predictability rather than one universal shortcut. The perturbation results further show that exact timestamp geometry is weakly necessary for competitive performance on several settings, while USHCN provides a clear counterexample where temporal arrangement matters.

\paragraph{Scope and limitations.}
The Anchor family is a diagnostic control, not a replacement for neural forecasting. Its variable-wise rules do not explicitly learn latent state, cross-variable dependencies, or regime dynamics, and the observed reversals at shorter lookbacks, longer horizons, and aggressive history thinning show where neural models regain advantages. These failures define the scope of the claim rather than weaken it: neural modeling is not broadly unnecessary, but its necessity is benchmark- and regime-dependent.

\paragraph{Implications for future IMTS evaluation.}
Future work should demonstrate improvements beyond strong controls for persistence, smoothing, recurrence, and non-test empirical selection, and should evaluate settings where irregular timing, asynchronous cross-variable effects, delayed dependencies, and distribution shifts are genuinely required for success. IMTS evaluation should distinguish between what an architecture \emph{can} represent and what a benchmark \emph{requires} it to represent. That distinction is necessary for making gains in model complexity scientifically interpretable.



\section{Conclusion}

This paper reframes strong non-neural baselines as diagnostic controls for IMTS forecasting. Across four benchmarks, the Anchor family shows that much reported performance can be matched by transparent history-derived rules under matched protocols. This does not imply neural models are unnecessary; the audits also reveal regimes where learned representations regain an advantage. Instead, we propose a stricter evaluation standard: future IMTS methods should demonstrate gains beyond persistence, smoothing, sparse repetition, recurrence, and non-test empirical selection, particularly where irregular temporal or cross-variable reasoning is required. 




\newpage

\section*{AI Use Statement}
In this work, ChatGPT, a Generative AI application, was used only to improve the clarity, fluency, and grammatical quality of the manuscript, as the authors are non-native English speakers. The authors provided the original words, technical content, arguments, structure, and key contexts, while AI assistance was limited to linking sentences, rephrasing passages, and polishing language. It  was not used to generate the research ideas, methodology, experiments, results, or scientific claims presented in this work. All AI-assisted text was reviewed and revised by the authors, who take full responsibility for the final content of the manuscript.


\bibliography{reference}
\bibliographystyle{iclr2026_conference}

\clearpage
\appendix


\section{Additional Anchor-family Analyses}
\label{app:additional-analyses}

This technical appendix expands the benchmark audit along three axes. First, it decomposes
which parts of the Anchor family are responsible for the standard-window
results. Second, it reports operational properties, including deterministic
CPU inference cost. Third, it studies whether the Anchor gains over pure
persistence are concentrated in particular volatility or observation-density
regimes. These analyses are not intended to introduce a second leaderboard;
they are included to make the diagnostic interpretation of the main results
auditable.

\paragraph{Candidate attribution.}
Tables~\ref{tab:candidate-removal} and~\ref{tab:anchor-family-ladder} show that
the standard-window Anchor results are not produced by one universal heuristic.
For P12, ERM selection over the shared candidate library accounts for the
improvement over NaiveAnchor: using rules only reverts to the NaiveAnchor
error, while removing trend has essentially no effect. MIMIC follows a similar
pattern, with ERMAnchor and \method{} nearly identical and substantially below
the rules-only condition in MSE. HumanActivity is different: the horizon-aware
EMA rule already reaches the best MSE, and removing EMA increases MSE from
0.0420 to 0.0483. USHCN is the clearest case where the structural components
matter: removing sparse and phase-sensitive candidates increases MSE by 0.0283
and 0.0616, respectively. Thus, all four standard settings are substantially
matched by history-derived controls, but the responsible low-complexity
structure differs by dataset.

\paragraph{Efficiency.}
Table~\ref{tab:anchor-efficiency} reports CPU inference measurements for the
Anchor variants. The absolute runtime differs by dataset size and number of
predicted coordinates, but all variants remain lightweight deterministic
procedures. The adaptive variants add modest overhead relative to NaiveAnchor
because they evaluate a richer rule set, yet \method{} still produces hundreds
to thousands of predictions per second on CPU. This matters for the audit
framing: the controls are not only non-neural, but also computationally simple
enough to serve as routine diagnostic baselines.

\paragraph{Temporal-information probes.}
Table~\ref{tab:temporal-tests} isolates the temporal information used by
\method{} itself. On P12 and HumanActivity, timestamp shuffling leaves the error
essentially unchanged, whereas keeping only time gaps causes large degradation.
This indicates that the strong Anchor performance in these settings can be
obtained from observed values and local history without strong dependence on
exact timestamp geometry. USHCN behaves differently: shuffling timestamps,
swapping history halves, or retaining only the last observation all substantially
increase error. This supports the main-text interpretation that AutoAnchor's
USHCN performance depends much more on temporal arrangement than its P12 or
HumanActivity performance.

\paragraph{Improvement over persistence.}
Table~\ref{tab:paired-anchor-tests} compares \method{} directly with
NaiveAnchor at the entity level. The gains over persistence are large and
statistically clear for P12, USHCN, and HumanActivity, and MIMIC shows a strong
MSE improvement even though the MAE interval includes zero. This paired
comparison is important because it separates the claim ``history is useful''
from the stronger claim made by the Anchor family: selecting among several
transparent history-derived structures gives a meaningfully stronger diagnostic
control than last-value persistence alone.

\paragraph{Volatility and density strata.}
Tables~\ref{tab:anchor-volatility-strata} and~\ref{tab:anchor-density-strata}
show where \method{} improves over NaiveAnchor. Across all four datasets, the
MSE reduction is positive in every volatility and density stratum. The largest
relative reductions often occur in more difficult strata: high-volatility
targets in P12, MIMIC, USHCN, and HumanActivity, and low- or high-density
targets in USHCN. These stratified results suggest that the Anchor improvements
are not confined to trivially stable or densely observed cases. Instead,
transparent history-derived rules continue to explain a substantial fraction of
error even when recent history is volatile or observation density changes.

In conclusion, the appendix reinforces the main diagnostic finding: strong performance across current IMTS benchmarks can often be explained by simple, transparent structures already present in the observed history. The specific mechanism varies across datasets, ranging from persistence and smoothing to recurrence and adaptive rule selection, but these controls remain lightweight and consistently improve over pure persistence across diverse volatility and observation-density regimes. At the same time, the temporal probes reveal meaningful dataset differences, with some benchmarks requiring little precise timestamp structure and others depending strongly on temporal arrangement.



\begin{table*}[t]
\caption{Candidate-removal ablations. Entries are global scaled MSE/MAE; deltas
are relative to full \method{} on the same dataset.}
\label{tab:candidate-removal}
\centering
\scriptsize
\resizebox{\textwidth}{!}{
\begin{tabular}{l|l|c|c|c|c}
\toprule
Dataset & Condition & MSE & $\Delta$MSE & MAE & $\Delta$MAE \\
\midrule
\multirow{5}{*}{P12} & full \method{} & 0.3020 & +0.0000 & 0.3585 & +0.0000 \\
& ERM only & 0.3020 & +0.0000 & 0.3585 & +0.0000 \\
& rules only & 0.3992 & +0.0972 & 0.3951 & +0.0366 \\
& minus EMA & 0.3157 & +0.0137 & 0.3683 & +0.0098 \\
& minus trend & 0.3019 & -0.0001 & 0.3585 & -0.0001 \\
\midrule
\multirow{5}{*}{USHCN} & full \method{} & 0.1565 & +0.0000 & 0.2187 & +0.0000 \\
& ERM only & 0.1662 & +0.0097 & 0.2634 & +0.0447 \\
& rules only & 0.1565 & +0.0000 & 0.2187 & +0.0000 \\
& minus sparse & 0.1848 & +0.0283 & 0.2719 & +0.0532 \\
& minus phase & 0.2181 & +0.0616 & 0.2605 & +0.0418 \\
\midrule
\multirow{5}{*}{HumanActivity} & full \method{} & 0.0420 & +0.0000 & 0.1148 & +0.0000 \\
& ERM only & 0.0428 & +0.0009 & 0.1131 & -0.0016 \\
& rules only & 0.0420 & +0.0000 & 0.1148 & +0.0000 \\
& minus EMA & 0.0483 & +0.0064 & 0.1193 & +0.0046 \\
& minus sparse/trend/phase & 0.0420 & +0.0000 & 0.1148 & +0.0000 \\
\midrule
\multirow{7}{*}{MIMIC} & full \method{} & 0.3785 & +0.0000 & 0.3657 & +0.0000 \\
& ERM only & 0.3782 & -0.0002 & 0.3658 & +0.0001 \\
& rules only & 0.4560 & +0.0775 & 0.3709 & +0.0053 \\
& minus EMA & 0.3878 & +0.0094 & 0.3684 & +0.0027 \\
& minus sparse & 0.3782 & -0.0002 & 0.3658 & +0.0001 \\
& minus trend & 0.3796 & +0.0011 & 0.3673 & +0.0016 \\
& minus phase & 0.3782 & -0.0003 & 0.3660 & +0.0003 \\
\bottomrule
\end{tabular}}
\end{table*}

\begin{table*}[t]
\centering
\caption{Anchor-family computational efficiency across datasets. \textbf{Bold} marks the best
value in each metric column, and \underline{underlining} marks the second best.
Lower is better for CPU Total and CPU ms/Pred.; higher is better for Pred./s.}
\label{tab:anchor-efficiency}
\scriptsize
\resizebox{\textwidth}{!}{
\begin{tabular}{l|c|c|c|c|c|c|c|c|c|c|c|c}
\toprule
Method
& \multicolumn{3}{|c|}{\begin{tabular}[c]{@{}c@{}}P12\end{tabular}}
& \multicolumn{3}{c|}{\begin{tabular}[c]{@{}c@{}}USHCN\end{tabular}}
& \multicolumn{3}{c|}{\begin{tabular}[c]{@{}c@{}}HumanActivity\end{tabular}}
& \multicolumn{3}{c}{\begin{tabular}[c]{@{}c@{}}MIMIC\end{tabular}} \\
\cmidrule(lr){2-4}
\cmidrule(lr){5-7}
\cmidrule(lr){8-10}
\cmidrule(lr){11-13}
& CPU Total (s) & ms/Pred. & Pred./s
& CPU Total (s) & ms/Pred. & Pred./s
& CPU Total (s) & ms/Pred. & Pred./s
& CPU Total (s) & ms/Pred. & Pred./s \\
\midrule
NaiveAnchor
& \second{9.1592} & \second{0.4245} & \second{2356.0}
& \second{0.3055} & \second{0.8301} & \second{1204.7}
& \second{1.5235} & \second{0.2578} & \second{3879.1}
& \best{18.6622} & \best{1.4117} & \best{708.4} \\

ExpoAnchor
& 9.3953 & 0.4354 & 2296.8
& \best{0.2843} & \best{0.7725} & \best{1294.4}
& 1.6863 & 0.2853 & 3504.7
& \second{18.7083} & \second{1.4151} & \second{706.6} \\

SparseAnchor
& \best{9.1545} & \best{0.4242} & \best{2357.2}
& 0.3069 & 0.8339 & 1199.2
& \best{1.4389} & \best{0.2435} & \best{4107.2}
& 19.1243 & 1.4466 & 691.3 \\

ERMAnchor
& 9.4662 & 0.4387 & 2279.6
& 0.3250 & 0.8832 & 1132.2
& 1.6848 & 0.2851 & 3507.8
& 19.1358 & 1.4475 & 690.9 \\

\method{}
& 9.6304 & 0.4463 & 2240.7
& 0.3414 & 0.9278 & 1077.8
& 1.6722 & 0.2829 & 3534.3
& 19.4374 & 1.4703 & 680.1 \\
\bottomrule
\end{tabular}
}
\end{table*}

\begin{table*}[t]
\caption{Anchor-family ablation ladder. Lower is better for MSE and MAE.
Bold marks the best Anchor-family value within each dataset; underline marks
the second best. Positive deltas indicate improvement over NaiveAnchor.}
\label{tab:anchor-family-ladder}
\centering
\scriptsize
\resizebox{\textwidth}{!}{
\begin{tabular}{l|l|c|c|c|c}
\toprule
Dataset & Method & MSE & $\Delta$ MSE vs Naive & MAE & $\Delta$ MAE vs Naive \\
\midrule
\multirow{5}{*}{\begin{tabular}[c]{@{}c@{}}P12\end{tabular}} & NaiveAnchor & 0.3992 & +0.0000 & 0.3951 & +0.0000 \\
 & ExpoAnchor & \second{0.3286} & +0.0706 & \second{0.3753} & +0.0198 \\
 & SparseAnchor & 0.3992 & +0.0000 & 0.3951 & +0.0000 \\
 & ERMAnchor & \best{0.3020} & +0.0972 & \best{0.3585} & +0.0366 \\
 & \method{} & \best{0.3020} & +0.0972 & \best{0.3585} & +0.0366 \\
\midrule
\multirow{5}{*}{\begin{tabular}[c]{@{}c@{}}USHCN\end{tabular}} & NaiveAnchor & 0.3640 & +0.0000 & 0.3291 & +0.0000 \\
 & ExpoAnchor & 0.5923 & -0.2284 & 0.5117 & -0.1826 \\
 & SparseAnchor & 0.2661 & +0.0978 & 0.2812 & +0.0479 \\
 & ERMAnchor & \second{0.1662} & +0.1978 & \second{0.2634} & +0.0657 \\
 & \method{} & \best{0.1565} & +0.2074 & \best{0.2187} & +0.1104 \\
\midrule
\multirow{5}{*}{\begin{tabular}[c]{@{}c@{}}HumanActivity\end{tabular}} & NaiveAnchor & 0.0605 & +0.0000 & 0.1225 & +0.0000 \\
 & ExpoAnchor & \best{0.0420} & +0.0185 & \second{0.1148} & +0.0078 \\
 & SparseAnchor & 0.0605 & +0.0000 & 0.1225 & +0.0000 \\
 & ERMAnchor & \second{0.0428} & +0.0177 & \best{0.1131} & +0.0094 \\
 & \method{} & \best{0.0420} & +0.0185 & \second{0.1148} & +0.0078 \\
\midrule
\multirow{5}{*}{\begin{tabular}[c]{@{}c@{}}MIMIC\end{tabular}} & NaiveAnchor & 0.4561 & +0.0000 & 0.3708 & +0.0000 \\
 & ExpoAnchor & 0.6431 & -0.1870 & 0.4974 & -0.1266 \\
 & SparseAnchor & 0.4560 & +0.0002 & 0.3709 & -0.0001 \\
 & ERMAnchor & \best{0.3782} & +0.0779 & \second{0.3658} & +0.0050 \\
 & \method{} & \second{0.3785} & +0.0777 & \best{0.3657} & +0.0051 \\
\bottomrule
\end{tabular}}
\end{table*}

\begin{table*}[t]
\caption{Temporal-information tests for \method{}. Values are global scaled
MSE/MAE; $\Delta$MSE is relative to the original history condition. Lower is
better.}
\label{tab:temporal-tests}
\centering
\scriptsize
\resizebox{\textwidth}{!}{
\begin{tabular}{l|l|c|c|c|c}
\toprule
Dataset & History condition & MSE & $\Delta$MSE & MAE & $\Delta$MAE \\
\midrule
\multirow{6}{*}{P12} & original & 0.3020 & +0.0000 & 0.3585 & +0.0000 \\
& shuffle timestamps & 0.3020 & -0.0000 & 0.3585 & -0.0001 \\
& swap halves & 0.4325 & +0.1305 & 0.4568 & +0.0983 \\
& drop early history & 0.3030 & +0.0010 & 0.3594 & +0.0008 \\
& keep last obs. & 0.3383 & +0.0363 & 0.3842 & +0.0256 \\
& keep time gaps only & 0.8017 & +0.4996 & 0.6930 & +0.3345 \\
\midrule
\multirow{5}{*}{USHCN} & original & 0.1565 & +0.0000 & 0.2187 & +0.0000 \\
& shuffle timestamps & 0.5159 & +0.3594 & 0.4301 & +0.2114 \\
& swap halves & 0.8154 & +0.6588 & 0.5435 & +0.3248 \\
& keep last obs. & 0.3640 & +0.2074 & 0.3291 & +0.1104 \\
& keep time gaps only & 0.4927 & +0.3362 & 0.4929 & +0.2742 \\
\midrule
\multirow{5}{*}{HumanActivity} & original & 0.0420 & +0.0000 & 0.1148 & +0.0000 \\
& shuffle timestamps & 0.0420 & +0.0000 & 0.1148 & +0.0000 \\
& swap halves & 0.1031 & +0.0611 & 0.1923 & +0.0775 \\
& keep last obs. & 0.0605 & +0.0185 & 0.1225 & +0.0078 \\
& keep time gaps only & 4.2491 & +4.2071 & 1.7014 & +1.5866 \\
\bottomrule
\end{tabular}}
\end{table*}

\begin{table*}[t]
\caption{Paired entity-level comparison of \method{} and NaiveAnchor. Deltas
are \method{} minus NaiveAnchor, so negative values favor \method{}.
Intervals use paired entity bootstrap; $p$-values use paired entity-level sign
flips. This comparison measures improvement over persistence within the Anchor
family and is not a paired test against neural models.}
\label{tab:paired-anchor-tests}
\centering
\scriptsize
\resizebox{\textwidth}{!}{
\begin{tabular}{l|l|c|c|c}
\toprule
Dataset & Metric & Delta (95\% CI) & $p$ & Entities \\
\midrule
\multirow{2}{*}{P12} & MAE & -0.0366 [-0.0436, -0.0299] & $<10^{-4}$ & 1,181 \\
& MSE & -0.0972 [-0.1125, -0.0820] & $<10^{-4}$ & 1,181 \\
\multirow{2}{*}{MIMIC} & MAE & -0.0051 [-0.0126, +0.0021] & 0.1714 & 1,534 \\
& MSE & -0.0777 [-0.1010, -0.0545] & $<10^{-4}$ & 1,534 \\
\multirow{2}{*}{USHCN} & MAE & -0.1104 [-0.1562, -0.0693] & $<10^{-4}$ & 112 \\
& MSE & -0.2074 [-0.3165, -0.1149] & $<10^{-4}$ & 112 \\
\multirow{2}{*}{HumanActivity} & MAE & -0.0078 [-0.0130, -0.0027] & 0.0041 & 218 \\
& MSE & -0.0185 [-0.0273, -0.0107] & $<10^{-4}$ & 218 \\
\bottomrule
\end{tabular}}
\end{table*}

\begin{table*}[t]
\caption{\method{} and NaiveAnchor errors stratified by recent-history
volatility. Positive MSE reduction favors \method{}; strata are
dataset-specific target-level tertiles.}
\label{tab:anchor-volatility-strata}
\centering
\scriptsize
\resizebox{\textwidth}{!}{
\begin{tabular}{l|l|c|c|c|c|c}
\toprule
Dataset & Stratum & Auto MAE & Naive MAE & Auto MSE & Naive MSE & MSE red. (\%) \\
\midrule
\multirow{3}{*}{P12} & low & \best{0.1878} & \second{0.1882} & \best{0.1280} & \second{0.1396} & +8.3 \\
& mid & \best{0.3860} & \second{0.4240} & \best{0.2851} & \second{0.3641} & +21.7 \\
& high & \best{0.5018} & \second{0.5733} & \best{0.4930} & \second{0.6939} & +29.0 \\
\multirow{3}{*}{MIMIC} & low & \second{0.2376} & \best{0.2136} & \best{0.2123} & \second{0.2170} & +2.1 \\
& mid & \best{0.3607} & \second{0.3646} & \best{0.2870} & \second{0.3290} & +12.8 \\
& high & \best{0.4987} & \second{0.5342} & \best{0.6361} & \second{0.8224} & +22.7 \\
\multirow{3}{*}{USHCN} & low & \best{0.0295} & \second{0.0320} & \best{0.0075} & \second{0.0122} & +38.3 \\
& mid & \best{0.3400} & \second{0.4622} & \best{0.2495} & \second{0.4717} & +47.1 \\
& high & \best{0.2875} & \second{0.4943} & \best{0.2133} & \second{0.6088} & +65.0 \\
\multirow{3}{*}{HumanActivity} & low & \best{0.0470} & \second{0.0535} & \best{0.0059} & \second{0.0072} & +18.0 \\
& mid & \best{0.1053} & \second{0.1125} & \best{0.0292} & \second{0.0337} & +13.4 \\
& high & \best{0.1921} & \second{0.2015} & \best{0.0908} & \second{0.1406} & +35.4 \\
\bottomrule
\end{tabular}}
\end{table*}

\begin{table*}[t]
\caption{\method{} and NaiveAnchor errors stratified by observation density.
Positive MSE reduction favors \method{}; strata are dataset-specific
target-level tertiles.}
\label{tab:anchor-density-strata}
\centering
\scriptsize
\resizebox{\textwidth}{!}{
\begin{tabular}{l|l|c|c|c|c|c}
\toprule
Dataset & Stratum & Auto MAE & Naive MAE & Auto MSE & Naive MSE & MSE red. (\%) \\
\midrule
\multirow{3}{*}{P12} & low & \best{0.3544} & \second{0.3753} & \best{0.2985} & \second{0.3778} & +21.0 \\
& mid & \best{0.3760} & \second{0.4170} & \best{0.3066} & \second{0.4138} & +25.9 \\
& high & \best{0.3452} & \second{0.3930} & \best{0.3009} & \second{0.4060} & +25.9 \\
\multirow{3}{*}{MIMIC} & low & \best{0.3642} & \second{0.3792} & \best{0.3207} & \second{0.4034} & +20.5 \\
& mid & \best{0.3822} & \second{0.3863} & \best{0.3786} & \second{0.4418} & +14.3 \\
& high & \second{0.3506} & \best{0.3470} & \best{0.4361} & \second{0.5232} & +16.6 \\
\multirow{3}{*}{USHCN} & low & \best{0.2565} & \second{0.3914} & \best{0.1781} & \second{0.4589} & +61.2 \\
& mid & \best{0.2236} & \second{0.2952} & \best{0.1832} & \second{0.3221} & +43.1 \\
& high & \best{0.1760} & \second{0.3005} & \best{0.1084} & \second{0.3105} & +65.1 \\
\multirow{3}{*}{HumanActivity} & low & \best{0.1340} & \second{0.1392} & \best{0.0510} & \second{0.0757} & +32.6 \\
& mid & \best{0.1111} & \second{0.1174} & \best{0.0391} & \second{0.0527} & +25.8 \\
& high & \best{0.0992} & \second{0.1110} & \best{0.0357} & \second{0.0531} & +32.7 \\
\bottomrule
\end{tabular}}
\end{table*}
\end{document}
